% !TEX root = ../main.tex
\corrigesChapitre{19}{Premier principe de la thermodynamique}
\corrigesSouples

\begin{blocLivre}{cor-19-1}
\begin{corrige}{1}{L'énergie interne}
\demarche{L'énergie interne concerne l'échelle microscopique ; le mouvement d'ensemble relève de l'énergie mécanique.}
\begin{questions}
\item L'énergie cinétique d'\res{\text{agitation}} des entités, les énergies potentielles d'\res{\text{interaction}} entre
entités et les énergies de \res{\text{liaison}} au sein des molécules et des noyaux.
\item Non : la température et l'état de la balle sont inchangés. C'est son \res{\text{énergie cinétique macroscopique}}
qui a augmenté.
\item Eau chauffée : énergie d'\res{\text{agitation}} ; glaçon qui fond : énergie d'\res{\text{interaction}} entre molécules ;
méthane qui brûle : énergies de \res{\text{liaison}} chimique.
\item Elle fait intervenir toutes les entités et toutes leurs interactions : on ne sait pas la calculer, et seules ses
\res{\text{variations}} se mesurent, à travers les échanges d'énergie.
\end{questions}
\end{corrige}
\end{blocLivre}
\begin{blocLivre}{cor-19-2}
\begin{corrige}{2}{Dans quel sens ?}
\demarche{Le transfert thermique va spontanément du corps chaud vers le corps froid ; il est positif pour le système qui
le reçoit.}
\begin{questions}
\item De la main vers la canette : $\res{Q_\text{canette}>0}$ et $\res{Q_\text{main}<0}$.
\item Le thé cède de l'énergie à l'air : $\res{Q<0}$. Le transfert cesse quand le thé atteint \qty{20}{\degreeCelsius}
(équilibre thermique).
\item Un \res{\text{travail}} : les forces de frottement se déplacent avec les mains.
\item Non : le froid n'est pas une énergie. C'est un \res{\text{transfert thermique}} qui sort de la pièce, plus chaude,
vers l'extérieur, plus froid.
\end{questions}
\end{corrige}
\end{blocLivre}
\begin{blocLivre}{cor-19-3}
\begin{corrige}{3}{Chauffer de l'eau ou de l'aluminium}
\demarche{$\Delta U=m\,c\,\Delta T$ ; un écart de \qty{45}{\degreeCelsius} vaut \qty{45}{K}.}
\begin{questions}
\item $\Delta U=2{,}0\times4180\times45=\res{\qty{3,76e5}{J}}$.
\item $\Delta t=\dfrac{\Delta U}{P}=\dfrac{\num{3.76e5}}{1500}=\res{\qty{251}{s}}$, soit 4~min~11~s.
\item $\Delta U=2{,}0\times897\times45=\res{\qty{8,07e4}{J}}$.
\item À masse et écart de température égaux, l'eau stocke \res{\text{4,7 fois plus}} d'énergie que l'aluminium : elle
transporte ou garde beaucoup d'énergie pour un faible écart de température.
\end{questions}
\end{corrige}
\end{blocLivre}
\begin{blocLivre}{cor-19-4}
\begin{corrige}{4}{La bouilloire}
\demarche{Système : l'eau ; premier principe $\Delta U=W+Q$, avec $W=P\,\Delta t$.}
\begin{questions}
\item $\Delta U=1{,}0\times4180\times80=\res{\qty{3,34e5}{J}}$.
\item $W=P\,\Delta t=2200\times180=\res{\qty{3,96e5}{J}}$.
\item $Q=\Delta U-W=\res{\qty{-6,2e4}{J}}$ : négatif, l'eau \res{\text{cède}} de l'énergie à la bouilloire et à l'air.
\item $\eta=\dfrac{\Delta U}{W}=\res{\qty{84}{\percent}}$.
\end{questions}
\end{corrige}
\end{blocLivre}
\begin{blocLivre}{cor-19-5}
\begin{corrige}{5}{Faire fondre un glaçon}
\demarche{Fusion à température constante : $Q=m\,L_\mathrm{f}$ ; chauffage de l'eau liquide : $Q=m\,c\,\Delta T$.}
\begin{questions}
\item $Q=0{,}030\times\num{3.34e5}=\res{\qty{1,00e4}{J}}$. La température \res{\text{reste égale à }\qty{0}{\degreeCelsius}}
pendant toute la fusion.
\item $Q'=0{,}030\times4180\times20=\res{\qty{2,5e3}{J}}$.
\item La fusion demande \res{\text{4 fois plus}} d'énergie que le chauffage de \qty{0}{\degreeCelsius} à \qty{20}{\degreeCelsius}.
\item $\Delta t=\dfrac{Q}{P}=\dfrac{\num{1.00e4}}{5{,}0}=\qty{2,0e3}{s}$, soit environ \res{\qty{33}{min}}.
\end{questions}
\end{corrige}
\end{blocLivre}
\begin{blocLivre}{cor-19-6}
\begin{corrige}{6}{Le freinage d'une voiture}
\demarche{$v=\qty{25}{m.s^{-1}}$ ; l'énergie mécanique perdue devient énergie interne des disques : $m_\mathrm{d}\,c\,\Delta T$.}
\begin{questions}
\item $E_\mathrm{c}=\frac12\times1200\times25^2=\res{\qty{3,75e5}{J}}$.
\item $\Delta T=\dfrac{E_\mathrm{c}}{4\times6{,}0\times460}=\res{\qty{34}{K}}$.
\item À vitesse constante, $E_\mathrm{c}$ ne change pas : les freins dissipent $m\,g\,h=1200\times9{,}81\times1000=\res{\qty{1,18e7}{J}}$.
\item $\Delta T=\res{\qty{1,07e3}{K}}$ : les disques seraient détruits. Ils cèdent heureusement de l'énergie à l'air, mais
s'échauffent dangereusement ; le \res{\text{frein moteur}} dissipe une partie de l'énergie dans le moteur.
\end{questions}
\end{corrige}
\end{blocLivre}
\begin{blocLivre}{cor-19-7}
\begin{corrige}{7}{Identifier un métal}
\demarche{Calorimètre : $\Delta U_\mathrm{total}=0$, soit $(m_\mathrm{e}c_\mathrm{e}+C_\mathrm{cal})(T_\mathrm{f}-T_\mathrm{e})+m\,c\,(T_\mathrm{f}-T_\mathrm{m})=0$.}
\begin{questions}
\item Le calorimètre empêche presque tout transfert thermique avec l'extérieur ($Q=0$) et aucun travail n'est échangé
($W=0$) : $\res{\Delta U=0}$.
\item $c=\dfrac{(m_\mathrm{e}c_\mathrm{e}+C_\mathrm{cal})(T_\mathrm{f}-T_\mathrm{e})}{m\,(T_\mathrm{m}-T_\mathrm{f})}=\dfrac{(836+40)\times3{,}4}{0{,}100\times76{,}6}=\res{\qty{389}{J.kg^{-1}.K^{-1}}}$.
\item C'est le \res{\text{cuivre}} : écart relatif de $\dfrac{389-385}{385}=\res{\qty{1,0}{\percent}}$.
\item Sans $C_\mathrm{cal}$ : $c=\qty{371}{J.kg^{-1}.K^{-1}}$, soit \qty{3,6}{\percent} d'écart. Il faut \res{\text{tenir compte}} de la
capacité thermique du calorimètre.
\end{questions}
\end{corrige}
\end{blocLivre}
\begin{blocLivre}{cor-19-8}
\begin{corrige}{8}{Calculer une température d'équilibre en Python}
\demarche{$\sum m_i\,c_i\,(T_\mathrm{f}-T_i)=0$ donne $T_\mathrm{f}=\dfrac{\sum m_i\,c_i\,T_i}{\sum m_i\,c_i}$.}
\begin{questions}
\item Chaque élément décrit un corps : nom, masse, capacité thermique massique, température initiale. \code{C} est la
\res{\text{capacité thermique totale}} $\sum m_i\,c_i$, \code{Tf} la \res{\text{température finale}}.
\item En développant la somme : $T_\mathrm{f}\sum m_i\,c_i=\sum m_i\,c_i\,T_i$. Les énergies reçues sont les variations d'énergie
interne des corps ; l'ensemble étant isolé, leur somme est \res{\text{nulle}}.
\item On ajoute \code{("calorimètre", 1.0, 40, 20.0)} : le produit $m\,c$ vaut alors \qty{40}{J.K^{-1}}. On obtient
$T_\mathrm{f}=\res{\qty{42,0}{\degreeCelsius}}$.
\item La glace \res{\text{fond}} à \qty{0}{\degreeCelsius} : une partie de l'énergie sert au changement d'état, $m\,L_\mathrm{f}$, que la
relation $\Delta U=m\,c\,\Delta T$ ne décrit pas.
\end{questions}
\end{corrige}
\end{blocLivre}
\begin{blocLivre}{cor-19-9}
\begin{corrige}{9}{Le chauffe-eau électrique}
\demarche{$\Delta U=m\,c\,\Delta T$ avec $m=\qty{200}{kg}$, puis $\Delta t=\Delta U/P$.}
\begin{questions}
\item $\Delta U=200\times4180\times45=\res{\qty{3,76e7}{J}}$, soit \res{\qty{10,45}{kWh}}.
\item $\Delta t=\dfrac{\num{3.762e7}}{2200}=\qty{1,71e4}{s}=\res{\qty{4,75}{h}}$ : le chauffage tient \res{\text{dans les heures creuses}}.
\item Coût : $10{,}45\times0{,}20\approx\res{2{,}09}$~€.
\item $\Delta T=\dfrac{1{,}5\times\num{3.6e6}}{200\times4180}=\res{\qty{6,5}{K}}$ : l'eau passe de \qty{60}{\degreeCelsius} à environ
\qty{53,5}{\degreeCelsius}.
\end{questions}
\end{corrige}
\end{blocLivre}
\begin{blocLivre}{cor-19-10}
\begin{corrige}{10}{Rafraîchir une boisson}
\demarche{Système isolé : $m_\mathrm{s}c\,(T_\mathrm{f}-25)+m_\mathrm{g}L_\mathrm{f}+m_\mathrm{g}c\,(T_\mathrm{f}-0)=0$.}
\begin{questions}
\item La glace \res{\text{fond}} à \qty{0}{\degreeCelsius}, puis l'eau de fonte \res{\text{se réchauffe}} jusqu'à la température finale.
\item $m_\mathrm{g}=\dfrac{m_\mathrm{s}c\times20}{L_\mathrm{f}+c\times5}=\dfrac{0{,}330\times4180\times20}{\num{3.34e5}+\num{2.09e4}}=\res{\qty{78}{g}}$ : un peu plus de
\res{\text{trois glaçons}}.
\item $T_\mathrm{f}=\dfrac{m_\mathrm{s}c\times25-m_\mathrm{g}L_\mathrm{f}}{(m_\mathrm{s}+m_\mathrm{g})\,c}=\res{\qty{5,6}{\degreeCelsius}}$.
\item $T_\mathrm{f}=\dfrac{330\times25}{405}=\res{\qty{20,4}{\degreeCelsius}}$ seulement : c'est surtout la \res{\text{fusion}} de la glace
qui refroidit la boisson.
\end{questions}
\end{corrige}
\end{blocLivre}
\begin{blocLivre}{cor-19-11}
\begin{corrige}{11}{La balle qui s'écrase}
\demarche{Énergie cinétique convertie en énergie interne : $m\,c\,\Delta T$, puis $m\,L$ pour la fusion.}
\begin{questions}
\item $E_\mathrm{c}=\frac12\times0{,}010\times350^2=\res{\qty{613}{J}}$, soit \res{\qty{6,13e4}{J.kg^{-1}}}.
\item $\Delta T=\dfrac{0{,}5\times\num{6.13e4}}{129}=\qty{237}{K}$ : la balle atteint \res{\qty{257}{\degreeCelsius}}.
\item Atteindre \qty{327}{\degreeCelsius} demande $129\times307=\qty{3,96e4}{J.kg^{-1}}$, moins que \qty{6,13e4}{J.kg^{-1}} : la
balle \res{\text{atteint sa température de fusion}}.
\item Il reste $\num{6.13e4}-\num{3.96e4}=\qty{2,16e4}{J.kg^{-1}}$ pour la fusion : $\dfrac{\num{2.16e4}}{\num{2.30e4}}$, soit
\res{\qty{94}{\percent}} du plomb fondrait.
\end{questions}
\end{corrige}
\end{blocLivre}
\begin{blocLivre}{cor-19-12}
\begin{corrige}{12}{Induction ou gaz ?}
\demarche{Énergie utile $m\,c\,\Delta T$ ; rendement $=$ énergie utile / énergie consommée ; proportions de la combustion.}
\begin{questions}
\item $\Delta U=2{,}0\times4180\times80=\res{\qty{6,69e5}{J}}$.
\item $W=2800\times280=\qty{7,84e5}{J}$, d'où $\eta=\res{\qty{85}{\percent}}$.
\item Énergie de combustion : $\dfrac{\num{6.69e5}}{0{,}40}=\qty{1,67}{MJ}$, soit $m(\ce{CH4})=\res{\qty{33,4}{g}}$ ($\qty{2,09}{mol}$) ; il se
forme autant de moles de \ce{CO2} : $m(\ce{CO2})=\res{\qty{92}{g}}$.
\item $W=\qty{0,218}{kWh}$, soit environ \res{\qty{11}{g}} de \ce{CO2} : en France, l'induction émet environ \res{\text{huit fois
moins}} de dioxyde de carbone que le gaz pour ce chauffage.
\end{questions}
\end{corrige}
\end{blocLivre}
\begin{blocLivre}{cor-19-13}
\begin{corrige}{13}{Stocker la chaleur du Soleil}
\demarche{Système incompressible : $\Delta U=m\,c\,\Delta T$ ; énergie électrique $=0{,}375\,\Delta U$.}
\begin{questions}
\item Les sels sont \res{\text{liquides}} entre 291 et \qty{384}{\degreeCelsius} : système incompressible sans changement d'état.
\item $\Delta U=\num{28.5e6}\times1500\times93=\res{\qty{3,98e12}{J}}$, soit \res{\qty{1,10e3}{MWh}}.
\item Énergie électrique : $0{,}375\times1105=\qty{414}{MWh}$, soit $\dfrac{414}{50}=\res{\qty{8,3}{h}}$, un peu plus que les
7,5~h annoncées : les réservoirs \res{\text{perdent de l'énergie}} par transfert thermique et les pompes consomment une
partie de la production.
\item La centrale produit \res{\text{après le coucher du soleil}}, au moment de la pointe de consommation du soir.
\end{questions}
\end{corrige}
\end{blocLivre}
\begin{blocLivre}{cor-19-14}
\begin{corrige}{14}{Pourquoi transpire-t-on ?}
\demarche{Énergie $=$ puissance $\times$ durée ; $\Delta T=Q/(m\,c)$ ; masse évaporée $=Q/L$.}
\begin{questions}
\item $Q=900\times3600=\res{\qty{3,24e6}{J}}$.
\item $\Delta T=\dfrac{\num{3.24e6}}{70\times3500}=\res{\qty{13,2}{K}}$ : une hausse mortelle.
\item Pour $\Delta T=\qty{2}{K}$ : $\Delta t=\dfrac{2\times70\times3500}{900}=\qty{544}{s}$, soit environ \res{\qty{9}{min}}.
\item $m=\dfrac{\num{3.24e6}}{\num{2.4e6}}=\res{\qty{1,35}{kg}}$ par heure : le coureur doit \res{\text{compenser cette perte d'eau}}
pour pouvoir continuer à transpirer et éviter la déshydratation.
\end{questions}
\end{corrige}
\end{blocLivre}
\begin{blocLivre}{cor-19-15}
\begin{corrige}{15}{La canette auto-chauffante}
\demarche{L'énergie libérée par la réaction vaut $n\times\qty{65,2}{kJ}$ ; 60~\% seulement chauffe le café.}
\begin{questions}
\item L'énergie de \res{\text{liaison}} (énergie chimique) du système diminue : elle est libérée au profit de l'agitation.
\item $Q=0{,}250\times4180\times40=\res{\qty{4,18e4}{J}}$.
\item Énergie libérée : $\dfrac{\num{4.18e4}}{0{,}60}=\qty{6,97e4}{J}$, d'où $n=\dfrac{\num{6.97e4}}{\num{65.2e3}}=\qty{1,07}{mol}$ et
$m=\res{\qty{60}{g}}$ d'oxyde de calcium.
\item $n(\ce{H2O})=n(\ce{CaO})$, soit $m=\res{\qty{19,2}{g}}$. On en met davantage pour que \res{\text{tout l'oxyde réagisse}} :
une partie de l'eau se vaporise.
\end{questions}
\end{corrige}
\end{blocLivre}
